\documentclass[11pt,a4paper]{article}
\usepackage[T1]{fontenc}
\usepackage{lmodern}
\usepackage{isabelle,isabellesym}
\usepackage{amsmath,amssymb}
\usepackage[margin=1in]{geometry}
\usepackage{pdfsetup}

\isabellestyle{it}
\isadroptag{theory}
\renewcommand{\isastyletext}{\normalsize\normalfont}
\renewcommand{\isastyletxt}{\normalfont}

\begin{document}

\title{The Lectures' Kraus Maps as Kraus Maps of the AFP}
\author{Dustin Bryant}
\maketitle

\begin{abstract}
The AFP entry \emph{Kraus Maps} (Unruh) defines Kraus maps on trace-class operators of arbitrary
Hilbert spaces; the formalised lectures on quantum computing and completely positive maps work
with matrices.  For finite-dimensional spaces we identify the two: a list of Kraus matrices gives a
Kraus family whose Kraus map is the lectures' $\rho\mapsto\sum_iW_i\rho W_i^*$, trace preserving
exactly when $\sum_iW_i^*W_i=I$.  With Choi's theorem of the lectures it follows that every
completely positive map of matrices, in Paulsen's definition, is a Kraus map of the AFP, and every
completely positive trace-preserving map a trace-preserving one.
\end{abstract}

\parindent 0pt\parskip 0.5ex

\input{Kraus_Maps_Bridge}

\end{document}
